% GENERATED FILE. Edit canonical lessons, then run npm run generate:books.
% canonical-source-hash: fnv1a64:0ee621a0b68858a6
% canonical-lessons: KA-C150-tole, KA-C150-oresu, KA-C150-gudisu, KA-C150-kattu, KA-C150-holi

\chapter{To wash, To wipe, To sweep, To tie, To sew}
\label{ch:ka-to-wash-to-wipe-to-sweep-to-tie-to-sew}
\hlchaptermodality{150}

\begin{chapteropening}
\textbf{By the end of this chapter:} \emph{I can say to wash, to wipe, to sweep, to tie and to sew.}

Complete the last lesson of chapter 150: i can say to wash, to wipe, to sweep, to tie and to sew.
\end{chapteropening}

\section[\texorpdfstring{\kn{ತೊಳೆ} (toḷe)}{ತೊಳೆ (toḷe)}]{\kn{ತೊಳೆ} (toḷe) — to wash}

\label{lesson:KA-C150-tole}



\begin{quote}
Before the new one: say the Kannada for to call, then the Kannada for to count.
\end{quote}

\subsection*{You'll want to know: \kn{ತೊಳೆ}}
\textbf{\kn{ತೊಳೆ}} — \emph{toḷe} — \textquotedblleft{}to wash\textquotedblright{}.

\begin{grammarlens}[title={What you've built}]
One more word for this chapter.
\end{grammarlens}

\subsection*{Your turn}
\begin{itemize}
\raggedright
  \item \emph{Say it:} \emph{toḷe}
  \item \emph{Say it:} \emph{toḷe}, once more
  \item \emph{Say it:} say \emph{toḷe}
  \item \emph{Recall:} say the Kannada for to call, then the Kannada for to count, then say \emph{toḷe} again
\end{itemize}

\begin{tcolorbox}[breakable,colback=teal!4,colframe=teal!35!black,title={Before you move on}]
What does \kn{ತೊಳೆ} mean? (\textquotedblleft{}To wash\textquotedblright{}.) Say it once more.
\end{tcolorbox}
\section[\texorpdfstring{\kn{ಒರೆಸು} (oresu)}{ಒರೆಸು (oresu)}]{\kn{ಒರೆಸು} (oresu) — to wipe}

\label{lesson:KA-C150-oresu}



\begin{quote}
Before the new one: say the Kannada for to count, then the Kannada for to wash.
\end{quote}

\subsection*{You'll want to know: \kn{ಒರೆಸು}}
\textbf{\kn{ಒರೆಸು}} — \emph{oresu} — \textquotedblleft{}to wipe\textquotedblright{}.

\begin{grammarlens}[title={What you've built}]
One more word for this chapter.
\end{grammarlens}

\subsection*{Your turn}
\begin{itemize}
\raggedright
  \item \emph{Say it:} \emph{oresu}
  \item \emph{Say it:} \emph{oresu}, once more
  \item \emph{Say it:} say \emph{oresu}
  \item \emph{Recall:} say the Kannada for to count, then the Kannada for to wash, then say \emph{oresu} again
\end{itemize}

\begin{tcolorbox}[breakable,colback=teal!4,colframe=teal!35!black,title={Before you move on}]
What does \kn{ಒರೆಸು} mean? (\textquotedblleft{}To wipe\textquotedblright{}.) Say it once more.
\end{tcolorbox}
\section[\texorpdfstring{\kn{ಗುಡಿಸು} (guḍisu)}{ಗುಡಿಸು (guḍisu)}]{\kn{ಗುಡಿಸು} (guḍisu) — to sweep}

\label{lesson:KA-C150-gudisu}



\begin{quote}
Before the new one: say the Kannada for to wash, then the Kannada for to wipe.
\end{quote}

\subsection*{You'll want to know: \kn{ಗುಡಿಸು}}
\textbf{\kn{ಗುಡಿಸು}} — \emph{guḍisu} — \textquotedblleft{}to sweep\textquotedblright{}.

\begin{grammarlens}[title={What you've built}]
One more word for this chapter.
\end{grammarlens}

\subsection*{Your turn}
\begin{itemize}
\raggedright
  \item \emph{Say it:} \emph{guḍisu}
  \item \emph{Say it:} \emph{guḍisu}, once more
  \item \emph{Say it:} say \emph{guḍisu}
  \item \emph{Recall:} say the Kannada for to wash, then the Kannada for to wipe, then say \emph{guḍisu} again
\end{itemize}

\begin{tcolorbox}[breakable,colback=teal!4,colframe=teal!35!black,title={Before you move on}]
What does \kn{ಗುಡಿಸು} mean? (\textquotedblleft{}To sweep\textquotedblright{}.) Say it once more.
\end{tcolorbox}
\section[\texorpdfstring{\kn{ಕಟ್ಟು} (kaṭṭu)}{ಕಟ್ಟು (kaṭṭu)}]{\kn{ಕಟ್ಟು} (kaṭṭu) — to tie, to build}

\label{lesson:KA-C150-kattu}



\begin{quote}
Before the new one: say the Kannada for to wipe, then the Kannada for to sweep.
\end{quote}

\subsection*{You'll want to know: \kn{ಕಟ್ಟು}}
\textbf{\kn{ಕಟ್ಟು}} — \emph{kaṭṭu} — \textquotedblleft{}to tie, to build\textquotedblright{}.

\begin{grammarlens}[title={What you've built}]
One more word for this chapter.
\end{grammarlens}

\subsection*{Your turn}
\begin{itemize}
\raggedright
  \item \emph{Say it:} \emph{kaṭṭu}
  \item \emph{Say it:} \emph{kaṭṭu}, once more
  \item \emph{Say it:} say \emph{kaṭṭu}
  \item \emph{Recall:} say the Kannada for to wipe, then the Kannada for to sweep, then say \emph{kaṭṭu} again
\end{itemize}

\begin{tcolorbox}[breakable,colback=teal!4,colframe=teal!35!black,title={Before you move on}]
What does \kn{ಕಟ್ಟು} mean? (\textquotedblleft{}To tie, to build\textquotedblright{}.) Say it once more.
\end{tcolorbox}
\section[\texorpdfstring{\kn{ಹೊಲಿ} (holi)}{ಹೊಲಿ (holi)}]{\kn{ಹೊಲಿ} (holi) — to sew}

\label{lesson:KA-C150-holi}



\begin{quote}
Before the new one: say the Kannada for to sweep, then the Kannada for to tie.
\end{quote}

\subsection*{You'll want to know: \kn{ಹೊಲಿ}}
\textbf{\kn{ಹೊಲಿ}} — \emph{holi} — \textquotedblleft{}to sew\textquotedblright{}.

\begin{grammarlens}[title={What you've built}]
One more word for this chapter.
\end{grammarlens}

\subsection*{Your turn}
\begin{itemize}
\raggedright
  \item \emph{Say it:} \emph{holi}
  \item \emph{Say it:} \emph{holi}, once more
  \item \emph{Say it:} say \emph{holi}
  \item \emph{Recall:} say the Kannada for to sweep, then the Kannada for to tie, then say \emph{holi} again
\end{itemize}

\begin{tcolorbox}[breakable,colback=teal!4,colframe=teal!35!black,title={Before you move on}]
What does \kn{ಹೊಲಿ} mean? (\textquotedblleft{}To sew\textquotedblright{}.) Say it once more.
\end{tcolorbox}
